\documentclass{uofa-eng-assignment}
\usepackage{caption}

\usepackage[spanish]{babel}
\usepackage[utf8]{inputenc}
\usepackage{amsmath,amssymb}
\usepackage{booktabs}
\usepackage{geometry}
\usepackage{hyperref}
\hypersetup{
    colorlinks = false,
    linkbordercolor = {white},
}
\captionsetup[table]{name=Tabla, labelsep=period}
\captionsetup[figure]{name=Figura, labelsep=period}

\newcommand{\hpl}{\hyperlink}
\newcommand{\hpt}{\hypertarget}
\newcommand*{\name}{Joan Sebastian Ramirez Sanchez  }
\newcommand*{\id}{202223948}
\newcommand*{\course}{Introducción a la teoría de control}
\newcommand*{\assignment}{Tarea 1 de 3}
\newcommand*{\dated}{\today}
\newcommand{\tf}{\textbf}

\begin{document}

\maketitle
\section*{6. Reactor Químico por Lotes (Batch)}

Un reactor donde ocurre la reacción exotérmica de Van de Vusse. Al operar por lotes, la concentración nominal del reactivo $C_{An}(t)$ decae aceleradamente disminuyendo la ganancia del control de alimentación.

\vspace{0.4cm}

\textbf{\large Modelo No Lineal}

\begin{equation}
\dot{C}_A
=
-k_1 C_A
-k_2 C_A^2
+\frac{1}{V}(C_{A0}-C_A)U(t)
\end{equation}

\vspace{0.3cm}

\textbf{\large Sistema Linealizado LTV}

\begin{equation}
A(t)=
\left[
-k_1
-2k_2 C_{An}(t)
-\frac{1}{V}u_n(t)
\right],
\qquad
B(t)=
\left[
\frac{1}{V}
\left(C_{A0}-C_{An}(t)\right)
\right].
\end{equation}

\vspace{0.3cm}

\textbf{Comportamiento interesante:}
La ganancia de control en $B(t)$ depende directamente de cuántos reactivos quedan libres. Al principio la acción del flujo $U(t)$ es fuerte, pero disminuye conforme el proceso llega a su fin.

\vspace{0.4cm}

\begin{center}
\begin{tabular}{llll}
\toprule
\textbf{Parámetro} & \textbf{Símbolo} & \textbf{Valor Sugerido} & \textbf{Unidad} \\
\midrule
Volumen del reactor & $V$ & 1,5 & m$^3$ \\
Concentración de la alimentación & $C_{A0}$ & 4,0 & mol/m$^3$ \\
Constante cinética lineal & $k_1$ & 0,05 & s$^{-1}$ \\
Constante cinética cuadrática & $k_2$ & 0,15 & m$^3$/(mol$\cdot$s) \\
\bottomrule
\end{tabular}
\end{center}

\vspace{0.4cm}

\textbf{\large Trayectoria Nominal de Prueba:}

Un decaimiento exponencial químico típico

\begin{equation}
C_{An}(t)=3.5\,e^{-0.03t}.
\end{equation}
\section*{Solución para algún forzante}

\end{document}
